\documentclass[11pt,letterpaper]{article}
\usepackage[utf8]{inputenc}
\usepackage[T1]{fontenc}
\usepackage[spanish,es-noshorthands,es-tabla]{babel}
\usepackage[margin=2cm,top=2.5cm,bottom=2.5cm]{geometry}
\usepackage{amsmath,amssymb}
\usepackage{enumitem}
\usepackage{fancyhdr}
\usepackage{listings}
\usepackage{xcolor}
\usepackage{tcolorbox}
\usepackage{booktabs}
\usepackage{array}
\raggedbottom
\setlength{\headheight}{14pt}
\let\vspaceoriginal\vspace
\renewcommand{\vspace}[1]{\par\vspaceoriginal{#1}\par}
\pagestyle{fancy}
\fancyhf{}
\fancyhead[L]{Basic C Without Tears}
\fancyhead[C]{\NombreDocumento}
\fancyhead[R]{Stack}
\fancyfoot[L]{Basic-C-Without-Tears}
\fancyfoot[R]{\thepage}
\renewcommand{\headrulewidth}{0.4pt}
\renewcommand{\footrulewidth}{0.4pt}
\lstdefinestyle{customc}{
  language=C,
  numbers=left,
  stepnumber=1,
  numbersep=8pt,
  numberstyle=\tiny\color{gray},
  basicstyle=\footnotesize\ttfamily,
  keywordstyle=\bfseries\color{blue!80!black},
  commentstyle=\itshape\color{green!50!black},
  stringstyle=\color{red!70!black},
  breaklines=true,
  breakatwhitespace=true,
  tabsize=4,
  showstringspaces=false,
  frame=none
}
\newtcolorbox{solucionbox}{
  before=\par,
  after=\par,
  colback=white,
  colframe=black,
  arc=0mm,
  boxrule=0.6pt,
  left=3mm,
  right=3mm,
  top=2mm,
  bottom=2mm
}
\newif\ifpauta
\ifdefined\PAUTA\pautatrue\else\pautafalse\fi

\newcommand{\NombreDocumento}{Examen 04 -- Variante B}
\begin{document}
\begin{center}
    {\LARGE \textbf{Basic C Without Tears}} \\[0.2cm]
    {\Large \textbf{Examen 04 -- Variante B}} \\[0.12cm]
    \small Fundamentos de C, matrices, validación y algoritmos simples de strings \\[0.08cm]
    \textbf{Autor:} Stack
\end{center}
\vspace{0.2cm}
\begin{enumerate}[label=\textbf{\arabic*.}]
\item Una matriz contiene resultados de dos equipos en dos rondas:
\begin{lstlisting}[style=customc, numbers=none]
int R[2][2] = {{3, 8}, {6, 4}};
\end{lstlisting}
\begin{enumerate}[label=(\alph*)]
\item Indique los valores de \texttt{R[0][1]} y \texttt{R[1][0]}.
\ifpauta\begin{solucionbox}\texttt{R[0][1]} vale 8 y \texttt{R[1][0]} vale 6.\end{solucionbox}\else\vspace{1.5cm}\fi
\item Calcule la suma de la diagonal principal y la diagonal secundaria.
\ifpauta\begin{solucionbox}Diagonal principal: $3+4=7$. Diagonal secundaria: $8+6=14$.\end{solucionbox}\else\vspace{2cm}\fi
\item Escriba código que calcule e imprima la suma de la diagonal principal.
\ifpauta\begin{solucionbox}\begin{lstlisting}[style=customc]
int suma = 0;
for (int i = 0; i < 2; i++) {
    suma += R[i][i];
}
printf("Diagonal principal: %d\n", suma);
\end{lstlisting}En la diagonal principal el índice de fila coincide con el de columna.\end{solucionbox}\else\vspace{3cm}\fi
\end{enumerate}
\newpage
\item Un número de unidades almacenado en \texttt{int unidades} es válido si está entre 1 y 500, inclusive.
\begin{enumerate}[label=(\alph*)]
\item Escriba la condición que comprueba la validez.
\ifpauta\begin{solucionbox}\texttt{unidades >= 1 \&\& unidades <= 500}.\end{solucionbox}\else\vspace{1.5cm}\fi
\item Escriba un fragmento \texttt{if-else} que imprima \texttt{Dentro de rango} o \texttt{Fuera de rango}.
\ifpauta\begin{solucionbox}\begin{lstlisting}[style=customc]
if (unidades >= 1 && unidades <= 500) {
    printf("Dentro de rango\n");
} else {
    printf("Fuera de rango\n");
}
\end{lstlisting}\end{solucionbox}\else\vspace{2cm}\fi
\item Clasifique \texttt{1}, \texttt{250}, \texttt{500} y \texttt{501}.
\ifpauta\begin{solucionbox}1, 250 y 500 son válidos; 501 está fuera del rango.\end{solucionbox}\else\vspace{2cm}\fi
\end{enumerate}
\newpage
\item Determine la salida exacta.
\begin{lstlisting}[style=customc]
int M[2][2] = {{1, 2}, {3, 4}};
for (int i = 0; i < 2; i++) {
    for (int j = 0; j < 2; j++) {
        if (M[i][j] % 2 == 0) {
            printf("%d ", M[i][j]);
        }
    }
}
printf("fin\n");
\end{lstlisting}
\ifpauta\begin{solucionbox}Se recorren las filas de izquierda a derecha. Solo 2 y 4 son pares.\par\textbf{Salida:} \texttt{2 4 fin}.\end{solucionbox}\else\vspace{3cm}\fi
\item Escriba una función que copie una cadena en otra en orden inverso. Suponga que el destino tiene espacio suficiente y que la cadena de entrada no supera los 99 caracteres.
\begin{center}\begin{tabular}{ll}\toprule\textbf{Entrada} & \textbf{Resultado}\\\midrule\texttt{"sol"} & \texttt{"los"}\\\texttt{"C"} & \texttt{"C"}\\\texttt{""} & \texttt{""}\\\bottomrule\end{tabular}\end{center}
\ifpauta\begin{solucionbox}\begin{lstlisting}[style=customc]
void invertir_copia(const char origen[], char destino[]) {
    int n = 0;
    while (origen[n] != '\0') {
        n++;
    }
    for (int i = 0; i < n; i++) {
        destino[i] = origen[n - 1 - i];
    }
    destino[n] = '\0';
}
\end{lstlisting}Primero se calcula la longitud; luego el índice fuente decrece implícitamente con la expresión $n-1-i$. Para longitud cero, no se ejecuta el \texttt{for} y se escribe el terminador en destino[0].\end{solucionbox}\else\vspace{3cm}\fi
\end{enumerate}
\end{document}
